\documentclass[11pt]{article}
\usepackage[utf8]{inputenc}
\usepackage{geometry}
\geometry{a4paper, margin=1in}
\usepackage{amsmath}
\usepackage{hyperref}
\usepackage{parskip}

\begin{document}

\begin{flushright}
\textbf{Date:} July 24, 2026 \\
\end{flushright}

\textbf{From:} \\
Manpreet Singh \\
Boston University \\
Boston, MA, 02215, USA \\
\href{mailto:manni@bu.edu}{manni@bu.edu} \\


\textbf{To:} \\
Guest Editors: Dr. Lydia Farina, Prof. Bernd Stahl, Dr. Helena Webb \\
Special Issue on ``Responsible use of AI in Art Creation and Archival Practice'' \\
\textit{AI \& Society: Journal of Knowledge, Culture and Communication} \\
Springer Nature \\

\vspace{.5em}

\noindent \textbf{SUBJECT:} Submission of Original Manuscript to Special Issue: \textit{Disentangling Algorithmic Bias from Archival Artifacts: A Controlled Audit of Vision-Language Model Valuation in Metropolitan Museum Archives}

\vspace{1em}

Dear Dr. Farina, Prof. Stahl, Dr. Webb, and Editorial Board,

We submit our original research manuscript titled \textbf{``Disentangling Algorithmic Bias from Archival Artifacts: A Controlled Audit of Vision-Language Model Valuation in Metropolitan Museum Archives''} for consideration in the Special Issue on \textbf{``Responsible use of AI in Art Creation and Archival Practice''} in \textit{AI \& Society: Journal of Knowledge, Culture and Communication}.

As vision-language models (VLMs) such as Contrastive Language-Image Pre-training (CLIP) are deployed across digital library indexing, museum collection search interfaces, and computational art history, auditing these systems for socio-cultural bias is vital for responsible AI governance in cultural heritage. A primary methodological challenge in archival AI auditing is distinguishing direct model evaluation disparities from confounding variables embedded within historical metadata—such as historical gender imbalances in physical artwork media (e.g., oil painting versus paper or textiles).

In this study, we establish a controlled quantitative audit framework evaluating zero-shot aesthetic valuation differential scores across OpenAI CLIP (trained on curated web data) and OpenCLIP (trained on uncurated LAION-2B web scrapes) using artwork metadata harvested from the Metropolitan Museum of Art Open Access API ($N=1,500$ harvested objects: $n=618$ unattributed/anonymous; $n=139$ missing/broken image URLs; final audited cohort $N=743$ attributed works, Male $n=534$, Female $n=209$).

Key contributions and empirical findings include:
\begin{itemize}
    \item \textbf{Unadjusted Model Auditing:} Baseline evaluations reveal no statistically significant difference in composite artistic value scores between male ($\mu = -0.0035$) and female ($\mu = -0.0067$) artists under OpenAI CLIP ($U = 59,307.00$, $p = 0.1829$, rank-biserial $r = -0.0628$) or OpenCLIP ($\mu_M = 0.0237, \mu_F = 0.0171$, $U = 59,867.00$, $p = 0.1224$, $r = -0.0728$).
    \item \textbf{Statistical Equivalence Testing (TOST):} Two One-Sided Tests confirm formal statistical score equivalence within a Cohen's $d = 0.30$ margin ($p_{\text{TOST}} = 0.0042$ for OpenAI CLIP; $p_{\text{TOST}} = 0.0024$ for OpenCLIP).
    \item \textbf{Multivariate Confound Control:} Ordinary Least Squares (OLS) regression controlling for artwork medium, creation era, and aspect ratio ($R^2 = 0.018, p = 0.512$ for CLIP; $R^2 = 0.017, p = 0.455$ for OpenCLIP) demonstrates artist gender exhibits no statistically significant main effect post-adjustment ($B = 0.0037, p = 0.202$ in CLIP; $B = 0.0065, p = 0.356$ in OpenCLIP), with low overall $R^2$ ($<1.8\%$) reflecting high residual embedding variance.
    \item \textbf{Responsible AI Archival Governance:} We propose institutional guidelines for digital archives, including medium-stratified z-score calibration ($z_{i,k} = \frac{S(x_i) - \mu_k}{\sigma_k}$) and multi-prompt sensitivity auditing to prevent automated cataloging systems from distorting public search rankings.
\end{itemize}

Given this Special Issue's focus on responsible AI frameworks for arts and cultural archives, our work provides an empirical methodology to prevent automated cataloging systems from amplifying historical representational gaps.

We confirm that:
\begin{enumerate}
    \item This work is original and is not under consideration, accepted, or published elsewhere.
    \item All authors have read and approved the manuscript for submission.
    \item There are no conflicts of interest to declare.
\end{enumerate}

Thank you for your time and consideration.

\vspace{1.5em}

\noindent Sincerely,

\vspace{1em}
\noindent Manpreet Singh \\
\noindent Boston University

\end{document}